We test how well coseismic landslide models trained on past earthquakes transfer to unseen ones, first across nine Pacific Rim events with open inventories, including Kaik\={o}ura 2016, then across 38 earthquakes worldwide, using only rapidly available inputs: Copernicus GLO-30 terrain, USGS ShakeMap ground motion and Newmark sliding-block displacement. Physics-based, logistic and gradient-boosting models rank hillslopes reasonably well in an unseen earthquake but mispredict its total amount of landsliding by up to three orders of magnitude, mostly as a shift in level rather than in shape. We therefore propose HiPhyBoost (hierarchical physics-anchored boosting). It combines (i) a Newmark-anchored, monotone-constrained boosted model with latent event offsets, (ii) an empirical-Bayes prior over offsets that yields zero-shot predictive intervals for the landslide total, and (iii) a quasi-likelihood Bayesian update of the offset from a few rapidly mapped image tiles that accounts for landslide clustering. Across 38 earthquakes, HiPhyBoost transferred better than physics-constrained boosting in 29 events (median AUC 0.795 against 0.763; mean gain 0.013), including 21 of 29 events not used in its development. It reached AUC 0.917 for Kaik\={o}ura without New Zealand training data. After one to ten mapped tiles it predicted the remaining landslide total within a median factor of 1.4--1.6, with far better calibrated intervals than a pixel-independent update. The event offsets summarise each earthquake's landslide propensity relative to generic slope mechanics but also reflect how inventories were mapped. Offsets tended to be higher after wetter-than-normal months over the shaken area, and conditioning the prior on covariates available at forecast time (magnitude, depth, mechanism, ShakeMap station count, antecedent wetness) narrowed zero-shot intervals more than threefold without loss of coverage, but did not improve point forecasts. For Aotearoa New Zealand, the results support issuing Alpine Fault or Hikurangi scenario forecasts as intervals and planning rapid tile mapping to narrow them.